\begin{examplebox}

	\textbf{\textcolor{CExample}{例 1（基础直接应用）}}

	\textbf{\textcolor{CExample}{题目：}}
	已知函数 $f(x)=x^2-2kx+1$ 对一切实数 $x$ 恒为正数，求实数 $k$ 的取值范围。

	\textbf{\textcolor{CExample}{解题步骤：}}
	\begin{description}[style=nextline, leftmargin=2.8em, labelsep=0.5em, itemsep=0.45em]
		\item[\textbf{第一步（识别开口）：}] 二次项系数 $a=1>0$，抛物线开口向上。
			\par\textbf{理由：}
			开口向上的二次函数在 $\mathbb{R}$ 上的最小值决定整体正负。
		\item[\textbf{第二步（计算判别式）：}]
			\[
				\begin{aligned}
					\Delta &= (-2k)^2 -4\cdot1\cdot1 \\
					&= 4k^2-4 .
			\end{aligned}
			\]
			\par\textbf{理由：}
			根据结论，只需 $\Delta<0$ 即可保证恒正。
		\item[\textbf{第三步（解不等式）：}] $4k^2-4<0 \Rightarrow k^2<1 \Rightarrow -1<k<1$。
			\par\textbf{理由：}
			求解一元二次不等式，得到 $k$ 的范围。
	\end{description}

	\textbf{\textcolor{CExample}{关键结论：}}
	\textbf{$k\in(-1,1)$}

	\medskip
	\textbf{\textcolor{CExample}{例 2（稍变形）}}

	\textbf{\textcolor{CExample}{题目：}}
	已知二次函数 $f(x)=kx^2-2x+k$（$k\neq0$）对一切实数 $x$ 恒为正，求 $k$ 的取值范围。

	\textbf{\textcolor{CExample}{解题步骤：}}
	\begin{description}[style=nextline, leftmargin=2.8em, labelsep=0.5em, itemsep=0.45em]
		\item[\textbf{第一步（判定开口）：}] 由恒正知 $k>0$（抛物线开口向上）。
			\par\textbf{理由：}
			若 $k<0$，开口向下，函数必有负值，与恒正矛盾。
		\item[\textbf{第二步（计算判别式）：}]
			\[
				\begin{aligned}
					\Delta &= (-2)^2 -4\cdot k\cdot k \\
					&= 4-4k^2 .
			\end{aligned}
			\]
			\par\textbf{理由：}
			充要条件要求 $\Delta<0$。
		\item[\textbf{第三步（解不等式合条件）：}] $4-4k^2<0 \Rightarrow k^2>1 \Rightarrow k>1$ 或 $k<-1$。结合 $k>0$，得 $k>1$。
			\par\textbf{理由：}
			取两个条件的公共部分。
	\end{description}

	\textbf{\textcolor{CExample}{关键结论：}}
	\textbf{$k\in(1,+\infty)$}

	\medskip
	\textbf{\textcolor{CExample}{例 3（综合应用）}}

	\textbf{\textcolor{CExample}{题目：}}
	已知二次函数 $f(x)=(m-1)x^2+(m-1)x+1$（$m\neq1$）对一切 $x\in\mathbb{R}$ 恒正，求 $m$ 的取值范围。

	\textbf{\textcolor{CExample}{解题步骤：}}
	\begin{description}[style=nextline, leftmargin=2.8em, labelsep=0.5em, itemsep=0.45em]
		\item[\textbf{第一步（判定开口）：}] 二次项系数 $m-1>0$，即 $m>1$。
			\par\textbf{理由：}
			开口向上是恒正的必要条件。
		\item[\textbf{第二步（计算判别式）：}]
			\[
				\begin{aligned}
					\Delta &= (m-1)^2 -4\cdot(m-1)\cdot1 \\
					&= (m-1)^2-4(m-1) \\
					&= (m-1)(m-5) .
			\end{aligned}
			\]
			\par\textbf{理由：}
			利用结论 $\Delta<0$。
		\item[\textbf{第三步（解不等式合并）：}] $(m-1)(m-5)<0 \Rightarrow 1<m<5$。结合 $m>1$，得 $1<m<5$。
			\par\textbf{理由：}
			开口条件已包含 $m>1$，取交集。
	\end{description}

	\textbf{\textcolor{CExample}{关键结论：}}
	\textbf{$m\in(1,5)$}

\end{examplebox}
